\begin{solution}{21}
    \begin{subsolution}
        设圆盘像到薄凸透镜的距离为 $v$。由题意知：$u = 20\mathrm{cm}$，$f = 10\mathrm{cm}$，代入透镜成像公式
        \begin{equation}
            \frac{1}{v} + \frac{1}{u} = \frac{1}{f}
            \label{eq:1.s21}
        \end{equation}
        得像距为
        \begin{equation}
            v = 20\mathrm{cm}
            \label{eq:2.s21}
        \end{equation}
        其横向放大率为
        \begin{equation}
            \beta = -\frac{v}{u} = -1
            \label{eq:3.s21}
        \end{equation}
        可知圆盘像在凸透镜右边 20cm，半径为 5cm，为圆盘状，圆盘与其像大小一样。
    \end{subsolution}
    \begin{subsolution}
        如下图所示，连接 A、B 两点，连线 AB 与光轴交点为 C 点，由两个相似三角形 $\triangle\mathrm{AOC}$ 与 $\triangle\mathrm{BB'C}$ 的关系可求得 C 点距离透镜为 $15\mathrm{cm}$。

        若将圆形光阑放置于凸透镜后方 $6\mathrm{cm}$ 处，此时圆形光阑在 C 点左侧。

        当圆形光阑半径逐渐减小时，均应有光线能通过圆形光阑在 B 点成像，因而圆盘像的形状及大小不变，而亮度变暗。

        此时不存在圆形光阑半径 $r_a$ 使得圆盘像大小的半径变为 (1) 中圆盘像大小的半径的一半。
    \end{subsolution}
    \begin{subsolution}
        若将圆形光阑移至凸透镜后方 $18\mathrm{cm}$ 处，此时圆形光阑在 C 点（距离透镜为 $15\mathrm{cm}$）的右侧。由下图所示，此时有：
        \begin{equation*}
            \mathrm{CB'} = \mathrm{BB'} = 5\mathrm{cm}, \quad \mathrm{R'B'} = 2\mathrm{cm},
        \end{equation*}
        利用两个相似三角形 $\triangle\mathrm{CRR'}$ 与 $\triangle\mathrm{CBB'}$ 的关系，得
        \begin{equation}
            r = \mathrm{RR'} = \frac{\mathrm{CR'}}{\mathrm{CB'}} \times \mathrm{BB'} = \frac{5 - 2}{5} \times 5\mathrm{cm} = 3\mathrm{cm}
            \label{eq:4.s21}
        \end{equation}
        可见当圆盘半径 $r = 3\mathrm{cm}$（光阑边缘与 AB 相交）时，圆盘刚好能成完整像，但其亮度变暗。

        若进一步减少光阑半径，圆盘像就会减小。当透镜上任何一点发出的光都无法透过光阑照在原先像的一半高度处时，圆盘像的半径就会减小为一半，如下图所示。此时光阑边缘与 AE 相交，AE 与光轴的交点为 D，由几何关系算得 D 与像的轴上距离为 $\frac{20}{7}\mathrm{cm}$。此时有
        \begin{equation*}
            \mathrm{DR'} = \frac{6}{7}\mathrm{cm}, \quad \mathrm{DE'} = \frac{20}{7}\mathrm{cm}, \quad \mathrm{EE'} = 2.5\mathrm{cm},
        \end{equation*}
        利用两个相似三角形 $\triangle\mathrm{DRR'}$ 与 $\triangle\mathrm{DEE'}$ 的关系，得
        \begin{equation}
            r_a = \mathrm{RR'} = \frac{\mathrm{DR'}}{\mathrm{DE'}} \times \mathrm{EE'} = \frac{20/7 - 2}{20/7} \times 2.5\mathrm{cm} = 0.75\mathrm{cm}
            \label{eq:5.s21}
        \end{equation}
        可见当圆形光阑半径 $r_a = 0.75\mathrm{cm}$，圆盘像大小的半径的确变为 (1) 中圆盘像大小的半径的一半。
    \end{subsolution}
    \begin{subsolution}
        只要圆形光阑放在 C 点（距离透镜为 $15\mathrm{cm}$）和光屏之间，圆盘像的大小便与圆形光阑半径有关。
    \end{subsolution}
    \begin{subsolution}
        若将图中的圆形光阑移至凸透镜前方 $6\mathrm{cm}$ 处，则当圆形光阑半径逐渐减小时，圆盘像的形状及大小不变，亮度变暗；

        同时不存在圆形光阑半径使得圆盘像大小的半径变为 (1) 中圆盘像大小的半径的一半。
    \end{subsolution}
\end{solution}